\documentclass[a4paper, 12pt]{article}
\usepackage[UTF8]{ctex} % 支持中文
\usepackage{geometry}
\geometry{left=2.0cm, right=2.0cm, top=2.5cm, bottom=2.5cm}
\usepackage{amsmath, amssymb}
\usepackage{booktabs} % 优化表格线条
\usepackage{longtable} % 长表格
\usepackage{array}
\usepackage{listings} % 代码环境
\usepackage{xcolor}
\usepackage{tikz} % 绘图（流程图）
\usepackage{multirow}
\usepackage{arydshln}
\usepackage{float} % 添加这一行
\usepackage{graphicx}
\usepackage{makecell}
\usetikzlibrary{shapes.geometric, arrows}
% 代码高亮设置
\lstset{
    language=Matlab,
    basicstyle=\small\ttfamily,
    keywordstyle=\color{blue},
    commentstyle=\color{green!60!black},
    stringstyle=\color{purple},
    numbers=left,
    numberstyle=\tiny\color{gray},
    stepnumber=1,
    frame=single,
    breaklines=true,
    captionpos=b,
    tabsize=4
}

\title{有限元法计算实例}
\author{}
\date{}

\begin{document}

\maketitle

\section{问题描述与网格剖分}

为了掌握有限元法解题的具体步骤与方法，必须结合具体问题进行分析。下面结合一个实际的数学例题，具体讨论 $\gamma$ 为常数时，有限元法解题的计算步骤及计算公式的应用。

如图 3.15 所示（注：此处指代原文档插图），有一正方形二维场区，域内电流密度 $J=0$，上、下边界处 $\frac{\partial A}{\partial n} = 0$，在左边界 $A=0$，右边界 $A=0.35$，介质是均匀的，磁阻系数为 $\gamma$，求场域内 $A$ 的分布。求解过程如下：

\textbf{第一步，进行场域剖分。}

\begin{enumerate}
    \item \textbf{节点坐标 $(x_n, y_n)$：}
    
    \begin{center}
    \begin{tabular}{ccc}
    1 (2.0, 0.0) & 2 (5.0, 0.0) & 3 (2.0, 3.0) \\
    4 (5.0, 3.0) & 5 (2.0, 5.0) & 6 (5.0, 5.0) \\
    7 (2.0, 7.0) & 8 (5.0, 7.0) & 9 (0.0, 0.0) \\
    10 (0.0, 3.0) & 11 (0.0, 5.0) & 12 (0.0, 7.0) \\
    13 (7.0, 0.0) & 14 (7.0, 3.0) & 15 (7.0, 5.0) \\
    16 (7.0, 7.0) & & \\
    \end{tabular}
    \end{center}

    \item \textbf{三角形单元编号及节点顺序 $(i, j, m)$：}
    
    \begin{center}
    \begin{tabular}{ccc}
    1 (15, 6, 4) & 2 (3, 4, 6) & 3 (1, 4, 3) \\
    4 (3, 5, 10) & 5 (3, 6, 5) & 6 (5, 6, 8) \\
    7 (6, 16, 8) & 8 (6, 15, 16) & 9 (4, 14, 15) \\
    10 (4, 13, 14) & 11 (4, 2, 13) & 12 (4, 1, 2) \\
    13 (3, 9, 1) & 14 (3, 10, 9) & 15 (5, 11, 10) \\
    16 (5, 12, 11) & 17 (5, 7, 12) & 18 (5, 8, 7) \\
    \end{tabular}
    \end{center}

    \item \textbf{强加边界条件节点编号及位函数值 $A$：}
    
    \begin{center}
    \begin{tabular}{cccc}
    (9, 0.0) & (10, 0.0) & (11, 0.0) & (12, 0.0) \\
    (13, 0.35) & (14, 0.35) & (15, 0.35) & (16, 0.35) \\
    \end{tabular}
    \end{center}
\end{enumerate}

\section{单元分析与 $K$ 矩阵计算}

\textbf{对每个单元进行单元分析，求出单元刚度系数矩阵 $K^{(e)}$，并同时进行整体合成。}

计算公式为：
\[
K_{st}^{(e)} = \frac{1}{4\Delta} (b_s b_t + c_s c_t)
\]

\begin{longtable}{|c|l|l|c|l|}
\caption{\textbf{ $K$ 矩阵元素计算列表}} \\
\hline
\textbf{单元} & $b_i = y_j - y_m$ & $c_i = x_m - x_j$ & \textbf{面积参数} & \textbf{系数矩阵元素} \\
\textbf{编号} & $b_j = y_m - y_i$ & $c_j = x_i - x_m$ & $4\Delta = 2(b_i c_j - b_j c_i)$ & $K_{st} = \frac{1}{4\Delta}(b_s b_t + c_s c_t)$ \\
 & $b_m = y_i - y_j$ & $c_m = x_j - x_i$ & & \\
\hline
\endfirsthead
\hline
\textbf{单元} & $b$ 计算 & $c$ 计算 & $4\Delta$ & \textbf{系数矩阵元素} \\
\hline
\endhead
\hline
\endfoot

% 单元 1
\multirow{3}{*}{1} & $b_{15}=5-3=2$ & $c_{15}=5-5=0$ & \multirow{3}{*}{$2(b_{15}c_6 - b_6 c_{15}) = 8$} & $K_{15,16} = \frac{1}{2} \quad K_{15,6} = -\frac{1}{2}$ \\
 & $b_6=3-5=-2$ & $c_6=7-5=2$ & & $K_{6,6}=1 \quad K_{15,4}=0$ \\
 & $b_4=5-5=0$ & $c_4=5-7=-2$ & & $K_{4,4}=\frac{1}{2} \quad K_{6,4} = -\frac{1}{2}$ \\
\hline

% 单元 2
\multirow{3}{*}{2} & $b_3=3-5=-2$ & $c_3=5-5=0$ & \multirow{3}{*}{12} & $K_{3,3}=\frac{4}{12} \quad K_{3,4} = -\frac{4}{12}$ \\
 & $b_4=5-3=2$ & $c_4=2-5=-3$ & & $K_{4,4}=\frac{19}{12} \quad K_{3,6}=0$ \\
 & $b_6=3-3=0$ & $c_6=5-2=3$ & & $K_{6,6}=\frac{21}{12} \quad K_{4,6} = -\frac{15}{12}$ \\
\hline

% 单元 3
\multirow{3}{*}{3} & $b_1=3-3=0$ & $c_1=2-5=-3$ & \multirow{3}{*}{18} & $K_{1,1}=\frac{1}{2} \quad K_{1,4}=0$ \\
 & $b_4=3$ & $c_4=2-2=0$ & & $K_{4,4}=\frac{25}{12} \quad K_{1,3} = -\frac{1}{2}$ \\
 & $b_3=0-3=-3$ & $c_3=5-2=3$ & & $K_{3,3}=\frac{8}{12} \quad K_{4,3} = -\frac{10}{12}$ \\
\hline

% 单元 4
\multirow{3}{*}{4} & $b_3=5-3=2$ & $c_3=0-2=-2$ & \multirow{3}{*}{8} & $K_{3,3}=\frac{28}{12} \quad K_{3,5} = -\frac{1}{2}$ \\
 & $b_5=3-3=0$ & $c_5=0$ & & $K_{5,5}=\frac{6}{12} \quad K_{3,10} = -\frac{1}{2}$ \\
 & $b_{10}=3-5=-2$ & $c_{10}=2-2=0$ & & $K_{10,10}=\frac{1}{2} \quad K_{5,10}=0$ \\
\hline

% 单元 5
\multirow{3}{*}{5} & $b_3=5-5=0$ & $c_3=2-5=-3$ & \multirow{3}{*}{12} & $K_{3,3}=\frac{37}{12} \quad K_{3,6}=0$ \\
 & $b_6=5-3=2$ & $c_6=2-2=0$ & & $K_{6,6}=\frac{30}{12} \quad K_{3,5} = -\frac{15}{12}$ \\
 & $b_5=3-5=-2$ & $c_5=5-2=3$ & & $K_{5,5}=\frac{19}{12} \quad K_{6,5} = -\frac{4}{12}$ \\
\hline

% 单元 6
\multirow{3}{*}{6} & $b_5=5-7=-2$ & $c_5=5-5=0$ & \multirow{3}{*}{12} & $K_{5,5}=\frac{23}{12} \quad K_{5,6} = -\frac{8}{12}$ \\
 & $b_6=7-5=2$ & $c_6=2-5=-3$ & & $K_{6,6}=\frac{38}{12} \quad K_{5,8}=0$ \\
 & $b_8=5-5=0$ & $c_8=5-2=3$ & & $K_{8,8}=\frac{9}{12} \quad K_{6,8} = -\frac{9}{12}$ \\
\hline

% 单元 7
\multirow{3}{*}{7} & $b_6=7-7=0$ & $c_6=5-7=-2$ & \multirow{3}{*}{8} & $K_{6,6}=\frac{44}{12} \quad K_{6,16}=0$ \\
 & $b_{16}=7-5=2$ & $c_{16}=5-5=0$ & & $K_{16,16}=\frac{1}{2} \quad K_{6,8} = -\frac{15}{12}$ \\
 & $b_8=5-7=-2$ & $c_8=7-5=2$ & & $K_{8,8}=\frac{21}{12} \quad K_{16,8} = -\frac{1}{2}$ \\
\hline

% 单元 8
\multirow{3}{*}{8} & $b_6=5-7=-2$ & $c_6=7-7=0$ & \multirow{3}{*}{8} & $K_{6,6}=\frac{50}{12} \quad K_{6,15} = -1$ \\
 & $b_{15}=7-5=2$ & $c_{15}=-7+5=-2$ & & $K_{15,15}=\frac{3}{2} \quad K_{6,16}=0$ \\
 & $b_{16}=5-5=0$ & $c_{16}=7-5=2$ & & $K_{16,16}=1 \quad K_{15,16} = -\frac{1}{2}$ \\
\hline

% 单元 9
\multirow{3}{*}{9} & $b_4=0-3=-3$ & $c_4=7-7=0$ & \multirow{3}{*}{8} & $K_{4,4}=\frac{31}{12} \quad K_{4,13} = -\frac{1}{2}$ \\
 & $b_{14}=5-3=2$ & $c_{14}=5-7=-2$ & & $K_{14,14}=1 \quad K_{4,15}=0$ \\
 & $b_{15}=3-3=0$ & $c_{15}=7-5=2$ & & $K_{15,15}=2 \quad K_{14,15} = -\frac{1}{2}$ \\
\hline

% 单元 10
\multirow{3}{*}{10} & $b_4=0-3=-3$ & $c_4=7-7=0$ & \multirow{3}{*}{12} & $K_{4,4}=\frac{40}{12} \quad K_{4,13} = 0$ \\
 & $b_2=3-3=0$ & $c_2=5-7=-2$ & & $K_{13,13}=\frac{4}{12} \quad K_{4,14}=-\frac{15}{12}$ \\
 & $b_{14}=3$ & $c_{14}=7-5=2$ & & $K_{14,14}=\frac{25}{12} \quad K_{13,14} = -\frac{4}{12}$ \\
\hline

% 单元 11
\multirow{3}{*}{11} & $b_4=0$ & $c_4=7-5=2$ & \multirow{3}{*}{12} & $K_{4,4}=\frac{44}{12} \quad K_{4,2} = -\frac{4}{12}$ \\
 & $b_2=0-3=-3$ & $c_2=5-7=-2$ & & $K_{2,2}=\frac{13}{12} \quad K_{4,13}=0$ \\
 & $b_{13}=3$ & $c_{13}=5-5=0$ & & $K_{13,13}=\frac{13}{12} \quad K_{2,13} = -\frac{9}{12}$ \\
\hline

% 单元 12
\multirow{3}{*}{12} & $b_4=0$ & $c_4=5-2=3$ & \multirow{3}{*}{18} & $K_{4,4}=\frac{50}{12} \quad K_{4,1}=0$ \\
 & $b_1=0-3=-3$ & $c_1=5-5=0$ & & $K_{1,1}=1 \quad K_{4,2} = -\frac{10}{12}$ \\
 & $b_2=3$ & $c_2=2-5=-3$ & & $K_{2,2}=\frac{25}{12} \quad K_{1,2} = -\frac{1}{2}$ \\
\hline

% 单元 13
\multirow{3}{*}{13} & $b_3=0$ & $c_3=2$ & \multirow{3}{*}{12} & $K_{3,3}=\frac{41}{12} \quad K_{3,9}=0$ \\
 & $b_9=-3$ & $c_9=3-2=1$ & & $K_{9,9}=\frac{9}{12} \quad K_{3,1} = -\frac{10}{12}$ \\
 & $b_1=3$ & $c_1=0-2=-2$ & & $K_{1,1}=\frac{25}{12} \quad K_{9,1} = -\frac{9}{12}$ \\
\hline

% 单元 14
\multirow{3}{*}{14} & $b_3=3$ & $c_3=0$ & \multirow{3}{*}{12} & $K_{3,3}=\frac{50}{12} \quad K_{3,10} = -\frac{15}{12}$ \\
 & $b_{10}=0-3=-3$ & $c_{10}=2$ & & $K_{10,10}=\frac{19}{12} \quad K_{3,9}=0$ \\
 & $b_9=0$ & $c_9=0-2=-2$ & & $K_{9,9}=\frac{13}{12} \quad K_{10,9} = -\frac{4}{12}$ \\
\hline

% 单元 15
\multirow{3}{*}{15} & $b_5=5-3=2$ & $c_5=0$ & \multirow{3}{*}{8} & $K_{5,5}=\frac{29}{12} \quad K_{5,11} = -\frac{1}{2}$ \\
 & $b_{11}=3-5=-2$ & $c_{11}=2$ & & $K_{11,11}=1 \quad K_{5,10}=0$ \\
 & $b_{10}=5-5=0$ & $c_{10}=0-2=-2$ & & $K_{10,10}=\frac{25}{12} \quad K_{11,10} = -\frac{1}{2}$ \\
\hline

% 单元 16
\multirow{3}{*}{16} & $b_5=7-5=2$ & $c_5=0$ & \multirow{3}{*}{8} & $K_{5,5}=\frac{35}{12} \quad K_{5,12}=0$ \\
 & $b_{12}=3-5=-2$ & $c_{12}=2$ & & $K_{12,12}=\frac{1}{2} \quad K_{5,11} = -1$ \\
 & $b_{11}=5-7=-2$ & $c_{11}=0-2=-2$ & & $K_{11,11}=2 \quad K_{12,11} = -\frac{1}{2}$ \\
\hline

% 单元 17
\multirow{3}{*}{17} & $b_5=7-7=0$ & $c_5=-2$ & \multirow{3}{*}{8} & $K_{5,5}=\frac{41}{12} \quad K_{5,7} = -\frac{1}{2}$ \\
 & $b_7=7-5=2$ & $c_7=2$ & & $K_{7,7}=-1 \quad K_{5,12}=0$ \\
 & $b_{12}=5-7=-2$ & $c_{12}=2-2=0$ & & $K_{12,12}=1 \quad K_{7,12} = -\frac{1}{2}$ \\
\hline

% 单元 18
\multirow{3}{*}{18} & $b_5=7-7=0$ & $c_5=2-5=-3$ & \multirow{3}{*}{12} & $K_{5,5}=\frac{50}{12} \quad K_{5,8}=0$ \\
 & $b_8=7-5=2$ & $c_8=2-2=0$ & & $K_{8,8}=\frac{25}{12} \quad K_{5,7} = -\frac{15}{12}$ \\
 & $b_7=5-7=-2$ & $c_7=5-2=3$ & & $K_{7,7}=\frac{25}{12} \quad K_{8,7} = -\frac{4}{12}$ \\
\hline
\end{longtable}

\section{整体刚度矩阵与方程组求解}

通过表 3.1 的计算结果可以列出所求的系数矩阵，它是一个 $16 \times 16$ 阶稀疏对称矩阵。按前面的讨论，$K$ 可以分成 4 块，因为节点有 16 个，其中包括 8 个已知节点。

\[
K = \begin{bmatrix} K_{11} & K_{12} \\ K_{12}^T & K_{22} \end{bmatrix}
\]

$K_{11}$ 是 $8 \times 8$ 的矩阵，如下虚线所分割：

\[
\resizebox{\textwidth}{!}{%
$
K = 
\left[
\begin{array}{cccccccc:cccccccc}
\frac{25}{12} & -\frac{1}{2} & -\frac{10}{12} & 0 & 0 & 0 & 0 & 0 & -\frac{9}{12} & 0 & 0 & 0 & 0 & 0 & 0 & 0 \\
-\frac{1}{2} & \frac{25}{12} & 0 & -\frac{10}{12} & 0 & 0 & 0 & 0 & 0 & 0 & 0 & 0 & -\frac{9}{12} & 0 & 0 & 0 \\
-\frac{10}{12} & 0 & \frac{50}{12} & -\frac{10}{12} & -\frac{15}{12} & 0 & 0 & 0 & 0 & -\frac{15}{12} & 0 & 0 & 0 & 0 & 0 & 0 \\
0 & -\frac{10}{12} & -\frac{10}{12} & \frac{50}{12} & 0 & -\frac{15}{12} & 0 & 0 & 0 & 0 & 0 & 0 & 0 & -\frac{15}{12} & 0 & 0 \\
0 & 0 & -\frac{15}{12} & 0 & \frac{50}{12} & -\frac{8}{12} & -\frac{15}{12} & 0 & 0 & 0 & -1 & 0 & 0 & 0 & 0 & 0 \\
0 & 0 & 0 & -\frac{15}{12} & -\frac{8}{12} & \frac{50}{12} & 0 & -\frac{15}{12} & 0 & 0 & 0 & 0 & 0 & 0 & -1 & 0 \\
0 & 0 & 0 & 0 & -\frac{15}{12} & 0 & \frac{25}{12} & -\frac{4}{12} & 0 & 0 & 0 & -\frac{1}{2} & 0 & 0 & 0 & 0 \\
0 & 0 & 0 & 0 & 0 & -\frac{15}{12} & -\frac{4}{12} & \frac{25}{12} & 0 & 0 & 0 & 0 & 0 & 0 & 0 & -\frac{1}{2} \\
\hdashline
-\frac{9}{12} & 0 & 0 & 0 & 0 & 0 & 0 & 0 & \frac{13}{12} & -\frac{4}{12} & 0 & 0 & 0 & 0 & 0 & 0 \\
0 & 0 & -\frac{15}{12} & 0 & 0 & 0 & 0 & 0 & -\frac{4}{12} & \frac{25}{12} & -\frac{1}{2} & 0 & 0 & 0 & 0 & 0 \\
0 & 0 & 0 & 0 & -1 & 0 & 0 & 0 & 0 & -\frac{1}{2} & 2 & -\frac{1}{2} & 0 & 0 & 0 & 0 \\
0 & 0 & 0 & 0 & 0 & 0 & -\frac{1}{2} & 0 & 0 & 0 & -\frac{1}{2} & 1 & 0 & 0 & 0 & 0 \\
0 & -\frac{9}{12} & 0 & 0 & 0 & 0 & 0 & 0 & 0 & 0 & 0 & 0 & \frac{13}{12} & -\frac{4}{12} & 0 & 0 \\
0 & 0 & 0 & -\frac{15}{12} & 0 & 0 & 0 & 0 & 0 & 0 & 0 & 0 & -\frac{4}{12} & \frac{25}{12} & 0 & 0 \\
0 & 0 & 0 & 0 & 0 & -1 & 0 & 0 & 0 & 0 & 0 & 0 & 0 & 0 & \frac{3}{2} & -\frac{1}{2} \\
0 & 0 & 0 & 0 & 0 & 0 & 0 & -\frac{1}{2} & 0 & 0 & 0 & 0 & 0 & 0 & -\frac{1}{2} & 1
\end{array}
\right]
$
}
\]
按式 (3.98) 有：
\[
\begin{cases}
K_{11}A_1 = R_1 - K_{12}d \\
A_{II} = d
\end{cases}
\]

已知 $d = [0.0 \quad 0.0 \quad 0.0 \quad 0.0 \quad 0.35 \quad 0.35 \quad 0.35 \quad 0.35]^T$，电流密度 $J=0$，式 (3.98) 在本例中的 $R_1$ 项为 [0]，即为：
\[
\begin{cases}
K_{11}A_1 = -K_{12}d \\
A_{II} = d
\end{cases}
\]

解此线性代数方程组得：
\[
A_1 = [A_1 \; A_2 \; A_3 \; A_4 \; A_5 \; A_6 \; A_7 \; A_8]^T
\]
\[
= [0.1 \; 0.25 \; 0.1 \; 0.25 \; 0.1 \; 0.25 \; 0.1 \; 0.25]^T
\]

通过分析上述表格的计算和方程组的求解可以看出，有限元法的各个环节在编制程序中易于标准化，对于同类问题，程序的通用有助于缩短解题周期。

在用有限元法计算时，首先要设计或搞清逻辑关系，然后才能编写或搞清程序，这就是要画出或读懂计算框图。对于上面讨论的实例，可以画出如图 3.16 所示的计算框图。

\section{计算流程图}

\begin{figure}[H]
    \centering
    \begin{tikzpicture}[node distance=1.5cm]
        \tikzstyle{startstop} = [rectangle, rounded corners, minimum width=3cm, minimum height=1cm,text centered, draw=black]
        \tikzstyle{process} = [rectangle, minimum width=3cm, minimum height=1cm, align=center, draw=black]
        \tikzstyle{arrow} = [thick,->,>=stealth]

        \node (start) [startstop] {开始};
        \node (pro1) [process, below of=start] {刚度矩阵元素和\\右端矢量赋值};
        \node (pro2) [process, below of=pro1] {计算各单元 $a, b, c, \Delta$};
        \node (pro3) [process, below of=pro2] {形成刚度矩阵和右端矢量};
        \node (pro4) [process, below of=pro3] {按一类边界修改方程组};
        \node (pro5) [process, below of=pro4] {解线性代数方程组};
        \node (pro6) [process, below of=pro5] {由 $A$ 求 $B$};
        \node (pro7) [process, below of=pro6] {打印输出};
        \node (stop) [startstop, below of=pro7] {结束};

        \draw [arrow] (start) -- (pro1);
        \draw [arrow] (pro1) -- (pro2);
        \draw [arrow] (pro2) -- (pro3);
        \draw [arrow] (pro3) -- (pro4);
        \draw [arrow] (pro4) -- (pro5);
        \draw [arrow] (pro5) -- (pro6);
        \draw [arrow] (pro6) -- (pro7);
        \draw [arrow] (pro7) -- (stop);
    \end{tikzpicture}
    \caption{图 3.16 计算框图}
\end{figure}

\section{源程序代码}

附录 6 图 3.16 计算框图源程序：

\begin{lstlisting}[language=Matlab, caption={有限元计算 MATLAB 源程序}]
% 首先给出节点坐标 (x,y)
x = [2,5,2,5,2,5,2,5,0,0,0,0,7,7,7,7];
y = [0,0,3,3,5,5,7,7,0,3,5,7,0,3,5,7];

% 对三角形单元的节点按逆时针顺序编号
ie = [15,3,1,3,3,5,6,6,4,4,4,4,3,3,5,5,5,5];
je = [6,4,4,5,6,6,16,15,14,13,2,1,9,10,11,12,7,8];
me = [4,6,3,10,5,8,8,16,15,14,13,2,1,9,10,11,12,7];

% 给出强加边界节点的位函数
d = [zeros(1,4), 0.35 * ones(1,4)];
k = zeros(16,16);

% 单元分析
for t = 1:18
    % 计算 b, c 和四倍 delta 的值
    b(ie(t)) = y(je(t)) - y(me(t));
    b(je(t)) = y(me(t)) - y(ie(t));
    b(me(t)) = y(ie(t)) - y(je(t));
    
    c(ie(t)) = x(me(t)) - x(je(t));
    c(je(t)) = x(ie(t)) - x(me(t));
    c(me(t)) = x(je(t)) - x(ie(t));
    
    sideta(t) = 2 * (b(ie(t)) * c(je(t)) - b(je(t)) * c(ie(t)));
    
    % 进行总体合成
    ff = [ie(t), je(t), me(t)];
    for f1 = 1:3
        p = ff(f1);
        for f2 = f1:3
            q = ff(f2);
            k(p,q) = k(p,q) + (b(p)*b(q) + c(p)*c(q)) ./ sideta(t);
            k(q,p) = k(p,q);
        end
    end
end

% 列出系数矩阵
k11 = k(1:8, 1:8);
k12 = k(1:8, 9:16);

% 求解方程
A = -k11 \ (k12 * d');

% 得到结果为
A = A'; % 转置显示
disp(A);
% A = 0.1000 0.2500 0.1000 0.2500 0.1000 0.2500 0.1000 0.2500
\end{lstlisting}

















\[
\begin{bmatrix}
\frac{25}{12} & -\frac{1}{2} & -\frac{10}{12} & 0 & 0 & 0 & 0 & 0 \\
-\frac{1}{2} & \frac{25}{12} & 0 & -\frac{10}{12} & 0 & 0 & 0 & 0 \\
-\frac{10}{12} & 0 & \frac{50}{12} & -\frac{10}{12} & -\frac{15}{12} & 0 & 0 & 0 \\
0 & -\frac{10}{12} & -\frac{10}{12} & \frac{50}{12} & 0 & -\frac{15}{12} & 0 & 0 \\
0 & 0 & -\frac{15}{12} & 0 & \frac{50}{12} & -\frac{8}{12} & -\frac{15}{12} & 0 \\
0 & 0 & 0 & -\frac{15}{12} & -\frac{8}{12} & \frac{50}{12} & 0 & -\frac{15}{12} \\
0 & 0 & 0 & 0 & -\frac{15}{12} & 0 & \frac{25}{12} & -\frac{4}{12} \\
0 & 0 & 0 & 0 & 0 & -\frac{15}{12} & -\frac{4}{12} & \frac{25}{12}
\end{bmatrix}
\begin{bmatrix}
A_1 \\ A_2 \\ A_3 \\ A_4 \\ A_5 \\ A_6 \\ A_7 \\ A_8
\end{bmatrix}\]
=
\[
-
\begin{bmatrix}
-\frac{9}{12} & 0 & 0 & 0 & 0 & 0 & 0 & 0 \\
0 & 0 & 0 & 0 & -\frac{9}{12} & 0 & 0 & 0 \\
0 & -\frac{15}{12} & 0 & 0 & 0 & 0 & 0 & 0 \\
0 & 0 & 0 & 0 & 0 & -\frac{15}{12} & 0 & 0 \\
0 & 0 & -1 & 0 & 0 & 0 & 0 & 0 \\
0 & 0 & 0 & 0 & 0 & 0 & -1 & 0 \\
0 & 0 & 0 & -\frac{1}{2} & 0 & 0 & 0 & 0 \\
0 & 0 & 0 & 0 & 0 & 0 & 0 & -\frac{1}{2}
\end{bmatrix}
\begin{bmatrix}
0 \\ 0 \\ 0 \\ 0 \\ 0.35 \\ 0.35 \\ 0.35 \\ 0.35
\end{bmatrix}
=
\begin{bmatrix}
0 \\
0.2625 \\
0 \\
0.4375 \\
0 \\
0.35 \\
0 \\
0.175
\end{bmatrix}
\]



\[
\begin{bmatrix}
\frac{25}{12} & -\frac{1}{2} & -\frac{10}{12} & 0 & 0 & 0 & 0 & 0 \\
-\frac{1}{2} & \frac{25}{12} & 0 & -\frac{10}{12} & 0 & 0 & 0 & 0 \\
-\frac{10}{12} & 0 & \frac{50}{12} & -\frac{10}{12} & -\frac{15}{12} & 0 & 0 & 0 \\
0 & -\frac{10}{12} & -\frac{10}{12} & \frac{50}{12} & 0 & -\frac{15}{12} & 0 & 0 \\
0 & 0 & -\frac{15}{12} & 0 & \frac{50}{12} & -\frac{8}{12} & -\frac{15}{12} & 0 \\
0 & 0 & 0 & -\frac{15}{12} & -\frac{8}{12} & \frac{50}{12} & 0 & -\frac{15}{12} \\
0 & 0 & 0 & 0 & -\frac{15}{12} & 0 & \frac{25}{12} & -\frac{4}{12} \\
0 & 0 & 0 & 0 & 0 & -\frac{15}{12} & -\frac{4}{12} & \frac{25}{12}
\end{bmatrix}
\]

\[
P = 
\begin{bmatrix}
0 & 1 & 0 & 0 & 0 & 0 & 0 & 0 \\ % P第1行 = 单位矩阵第2行 → 取A的第2行作为结果第1行
0 & 0 & 0 & 1 & 0 & 0 & 0 & 0 \\ % P第2行 = 单位矩阵第4行 → 取A的第4行作为结果第2行
0 & 0 & 0 & 0 & 0 & 1 & 0 & 0 \\ % P第3行 = 单位矩阵第6行 → 取A的第6行作为结果第3行
0 & 0 & 0 & 0 & 0 & 0 & 0 & 1 \\ % P第4行 = 单位矩阵第8行 → 取A的第8行作为结果第4行
1 & 0 & 0 & 0 & 0 & 0 & 0 & 0 \\ % P第5行 = 单位矩阵第1行 → 取A的第1行作为结果第5行
0 & 0 & 1 & 0 & 0 & 0 & 0 & 0 \\ % P第6行 = 单位矩阵第3行 → 取A的第3行作为结果第6行
0 & 0 & 0 & 0 & 1 & 0 & 0 & 0 \\ % P第7行 = 单位矩阵第5行 → 取A的第5行作为结果第7行
0 & 0 & 0 & 0 & 0 & 0 & 1 & 0 \\ % P第8行 = 单位矩阵第7行 → 取A的第7行作为结果第8行
\end{bmatrix}
\]


\[
M=P K = 
\begin{bmatrix}
% 新第1行 = 原矩阵A的第2行
-\frac{1}{2} & \frac{25}{12} & 0 & -\frac{10}{12} & 0 & 0 & 0 & 0 \\
% 新第2行 = 原矩阵A的第4行
0 & -\frac{10}{12} & -\frac{10}{12} & \frac{50}{12} & 0 & -\frac{15}{12} & 0 & 0 \\
% 新第3行 = 原矩阵A的第6行
0 & 0 & 0 & -\frac{15}{12} & -\frac{8}{12} & \frac{50}{12} & 0 & -\frac{15}{12} \\
% 新第4行 = 原矩阵A的第8行
0 & 0 & 0 & 0 & 0 & -\frac{15}{12} & -\frac{4}{12} & \frac{25}{12} \\
% 新第5行 = 原矩阵A的第1行
\frac{25}{12} & -\frac{1}{2} & -\frac{10}{12} & 0 & 0 & 0 & 0 & 0 \\
% 新第6行 = 原矩阵A的第3行
-\frac{10}{12} & 0 & \frac{50}{12} & -\frac{10}{12} & -\frac{15}{12} & 0 & 0 & 0 \\
% 新第7行 = 原矩阵A的第5行
0 & 0 & -\frac{15}{12} & 0 & \frac{50}{12} & -\frac{8}{12} & -\frac{15}{12} & 0 \\
% 新第8行 = 原矩阵A的第7行
0 & 0 & 0 & 0 & -\frac{15}{12} & 0 & \frac{25}{12} & -\frac{4}{12}
\end{bmatrix}
\]


\[
P\mathbf{b} = 
\begin{bmatrix}
0.2625 \\ 0.4375 \\ 0.35 \\ 0.175 \\ 0 \\ 0 \\ 0 \\ 0
\end{bmatrix}
\]
\[A=\widetilde{M}
\]
\[
b=\widetilde{b}
\]
\[A x=b
\]
\[
\widetilde{M} =
\begin{bmatrix}
\mathbf{0}_{8\times8} & M \\
M^{\mathrm{T}} & \mathbf{0}_{8\times8}
\end{bmatrix}
\]
\[
b=
\begin{bmatrix}
0.2625 \\ 0.4375 \\ 0.35 \\ 0.175 \\ 0 \\ 0 \\ 0 \\ 0\\
0 \\ 0 \\ 0 \\ 0 \\ 0 \\ 0 \\ 0 \\ 0
\end{bmatrix}
\]
\[
A = 
\begin{array}{cccccccccccccccc}
0 & 0 & 0 & 0 & 0 & 0 & 0 & 0 & -\frac{1}{2} & \frac{25}{12} & 0 & -\frac{10}{12} & 0 & 0 & 0 & 0 \\
0 & 0 & 0 & 0 & 0 & 0 & 0 & 0 & 0 & -\frac{10}{12} & -\frac{10}{12} & \frac{50}{12} & 0 & -\frac{15}{12} & 0 & 0 \\
0 & 0 & 0 & 0 & 0 & 0 & 0 & 0 & 0 & 0 & 0 & -\frac{15}{12} & -\frac{8}{12} & \frac{50}{12} & 0 & -\frac{15}{12} \\
0 & 0 & 0 & 0 & 0 & 0 & 0 & 0 & 0 & 0 & 0 & 0 & 0 & -\frac{15}{12} & -\frac{4}{12} & \frac{25}{12} \\
0 & 0 & 0 & 0 & 0 & 0 & 0 & 0 & \frac{25}{12} & -\frac{1}{2} & -\frac{10}{12} & 0 & 0 & 0 & 0 & 0 \\
0 & 0 & 0 & 0 & 0 & 0 & 0 & 0 & -\frac{10}{12} & 0 & \frac{50}{12} & -\frac{10}{12} & -\frac{15}{12} & 0 & 0 & 0 \\
0 & 0 & 0 & 0 & 0 & 0 & 0 & 0 & 0 & 0 & -\frac{15}{12} & 0 & \frac{50}{12} & -\frac{8}{12} & -\frac{15}{12} & 0 \\
0 & 0 & 0 & 0 & 0 & 0 & 0 & 0 & 0 & 0 & 0 & 0 & -\frac{15}{12} & 0 & \frac{25}{12} & -\frac{4}{12} \\
-\frac{1}{2} & 0 & 0 & 0 & \frac{25}{12} & -\frac{10}{12} & 0 & 0 & 0 & 0 & 0 & 0 & 0 & 0 & 0 & 0 \\
\frac{25}{12} & -\frac{10}{12} & 0 & 0 & -\frac{1}{2} & 0 & 0 & 0 & 0 & 0 & 0 & 0 & 0 & 0 & 0 & 0 \\
0 & -\frac{10}{12} & 0 & 0 & -\frac{10}{12} & \frac{50}{12} & -\frac{15}{12} & 0 & 0 & 0 & 0 & 0 & 0 & 0 & 0 & 0 \\
-\frac{10}{12} & \frac{50}{12} & -\frac{15}{12} & 0 & 0 & -\frac{10}{12} & 0 & 0 & 0 & 0 & 0 & 0 & 0 & 0 & 0 & 0 \\
0 & 0 & -\frac{8}{12} & 0 & 0 & -\frac{15}{12} & \frac{50}{12} & -\frac{15}{12} & 0 & 0 & 0 & 0 & 0 & 0 & 0 & 0 \\

0 & -\frac{15}{12} & \frac{50}{12} & -\frac{15}{12} & 0 & 0 & -\frac{8}{12} & 0 & 0 & 0 & 0 & 0 & 0 & 0 & 0 & 0 \\
0 & 0 & 0 & -\frac{4}{12} & 0 & 0 & -\frac{15}{12} & \frac{25}{12} & 0 & 0 & 0 & 0 & 0 & 0 & 0 & 0 \\
0 & 0 & -\frac{15}{12} & \frac{25}{12} & 0 & 0 & 0 & -\frac{4}{12} & 0 & 0 & 0 & 0 & 0 & 0 & 0 & 0 \\
\end{array}
\]
\[
\widetilde{M} x=\widetilde{b}
\]
\[
\begin{array}{cccccccccccccccc}
7 & 0 & 0 & 0 & 0 & 0 & 0 & 0 & -\frac{1}{2} & \frac{25}{12} & 0 & -\frac{10}{12} & 0 & 0 & 0 & 0 \\
0 & 7 & 0 & 0 & 0 & 0 & 0 & 0 & 0 & -\frac{10}{12} & -\frac{10}{12} & \frac{50}{12} & 0 & -\frac{15}{12} & 0 & 0 \\
0 & 0 & 7 & 0 & 0 & 0 & 0 & 0 & 0 & 0 & 0 & -\frac{15}{12} & -\frac{8}{12} & \frac{50}{12} & 0 & -\frac{15}{12} \\
0 & 0 & 0 & 7 & 0 & 0 & 0 & 0 & 0 & 0 & 0 & 0 & 0 & -\frac{15}{12} & -\frac{4}{12} & \frac{25}{12} \\
0 & 0 & 0 & 0 & 7 & 0 & 0 & 0 & \frac{25}{12} & -\frac{1}{2} & -\frac{10}{12} & 0 & 0 & 0 & 0 & 0 \\
0 & 0 & 0 & 0 & 0 & 7 & 0 & 0 & -\frac{10}{12} & 0 & \frac{50}{12} & -\frac{10}{12} & -\frac{15}{12} & 0 & 0 & 0 \\
0 & 0 & 0 & 0 & 0 & 0 & 7 & 0 & 0 & 0 & -\frac{15}{12} & 0 & \frac{50}{12} & -\frac{8}{12} & -\frac{15}{12} & 0 \\
0 & 0 & 0 & 0 & 0 & 0 & 0 & 7 & 0 & 0 & 0 & 0 & -\frac{15}{12} & 0 & \frac{25}{12} & -\frac{4}{12} \\
-\frac{1}{2} & 0 & 0 & 0 & \frac{25}{12} & -\frac{10}{12} & 0 & 0 & 7 & 0 & 0 & 0 & 0 & 0 & 0 & 0 \\
\frac{25}{12} & -\frac{10}{12} & 0 & 0 & -\frac{1}{2} & 0 & 0 & 0 & 0 & 7 & 0 & 0 & 0 & 0 & 0 & 0 \\
0 & -\frac{10}{12} & 0 & 0 & -\frac{10}{12} & \frac{50}{12} & -\frac{15}{12} & 0 & 0 & 0 & 7 & 0 & 0 & 0 & 0 & 0 \\
-\frac{10}{12} & \frac{50}{12} & -\frac{15}{12} & 0 & 0 & -\frac{10}{12} & 0 & 0 & 0 & 0 & 0 & 7 & 0 & 0 & 0 & 0 \\
0 & 0 & -\frac{8}{12} & 0 & 0 & -\frac{15}{12} & \frac{50}{12} & -\frac{15}{12} & 0 & 0 & 0 & 0 & 7 & 0 & 0 & 0 \\

0 & -\frac{15}{12} & \frac{50}{12} & -\frac{15}{12} & 0 & 0 & -\frac{8}{12} & 0 & 0 & 0 & 0 & 0 & 0 & 7 & 0 & 0 \\
0 & 0 & 0 & -\frac{4}{12} & 0 & 0 & -\frac{15}{12} & \frac{25}{12} & 0 & 0 & 0 & 0 & 0 & 0 & 7 & 0 \\
0 & 0 & -\frac{15}{12} & \frac{25}{12} & 0 & 0 & 0 & -\frac{4}{12} & 0 & 0 & 0 & 0 & 0 & 0 & 0 & 7 \\
\end{array}
\]
最大正特征值为
\end{document}